\subsection{A complete noncommutative Galois--pseudo-unitary reduction}
\label{subsec:frontier-nc-galois-pu}

\begin{proposition}
Every complex fusion categorification of $R8\_3\_721$ or $R8\_3\_1154$
admits a spherical structure and has a pseudo-unitary Galois conjugate.
\end{proposition}
\begin{proof}
For each ring, the complexified fusion algebra is
$\mathbb C^4\oplus M_2(\mathbb C)$. The common kernel of all commutators
in the regular representation consists exactly of the scalar blocks.
The four characters obtained there exhaust all possible dimension
characters; a matrix block has no one-dimensional representation.
Its formal codegree is nevertheless included: if its regular Casimir
eigenvalue is $\lambda$, with multiplicity four, its formal codegree is
$\lambda/2$. Exact multiplication identities and the complete Casimir
factorization verify this accounting. For each original ring, every scalar-character orbit
except the FP orbit is excluded by a zero dimension or formal-codegree
norm divisibility.

For an arbitrary categorification take its canonical sphericalization
and choose dual-compatible lifts. The signed quotient by the sign object
$\varepsilon=-1$ satisfies the full parity intervals
$|T_{ij}^k|\leq N_{ij}^k$ and $T_{ij}^k\equiv N_{ij}^k\pmod2$,
cyclic trace symmetry, dual reversal, and every ordered associativity
equation. We impose neither commutativity nor a modulo-four indicator
normalization. Complete enumeration gives $64$ signed tensors for
$R8\_3\_721$, all gauge-equivalent to the original tensor, and $288$
for $R8\_3\_1154$. The latter comprise one $64$-member original gauge
class and five other classes of sizes $64,64,32,32,32$.
All scalar-character factors of the five other classes are excluded by
zero dimensions or exact norm divisibility. Both sign sectors' formal
codegrees, including matrix blocks, are used: sphericalization has
dimension $2\delta$ and formal codegrees $2f$, so central induction
requires $\delta/f$ to be integral.

Thus every possible signed quotient is gauge-original. Switching the
lifts by its dual-compatible gauge makes the signed tensor equal to
$N$. The negative lift multiplicities are then zero, so the chosen
lifts form a spherical fusion subcategory whose forgetful functor to
the original category is a tensor equivalence. Its dimension character
lies in the FP Galois orbit. A number-field model and Galois conjugation
therefore give a spherical category with positive FP dimensions, hence
a pseudo-unitary category with the same fusion rules.
\end{proof}

The source-bound replay is
\texttt{central\_integrality/general\_galois\_pu\_gate.py} in the frontier
evidence directory; its primary calculation incorporates the independent
unrestricted signed enumeration. The earlier normalized enumeration is
retained only as a cross-check. Consequently the central-support, partial
modular, integral Gram, joint-rotation, and low Verlinde constraints for
these two rings are necessary for general complex existence after Galois
transport. Passing these constraints supplies neither a model nor
unitarity; both rings remain unresolved by this proposition.
